% Background and Motivation: concepts, limitations, evidence, and challenges.
\section{Background and Motivation}
\label{sec:background-motivation}

\subsection{Background}
\label{sec:motivation-background}

\paragraph{Multi-agent execution.}
Many LLM applications rely on multi-agent collaboration to complete
complex tasks. These applications divide a user task into subtasks and
assign them to several agents for parallel execution. Each agent maintains
its own conversation, advancing its task through LLM requests and tool
invocations: the model generates responses or initiates tool calls, and
the returned results become input to subsequent
requests~\cite{wu2023autogen}. A \emph{request} is a single model call;
an \emph{agent} can initiate multiple requests to complete a subtask;
a \emph{workflow} connects these tasks through serial dependencies,
parallel branches and joins. The application orchestrates the workflow,
while the inference service executes its model requests.

Figure~\ref{fig:agent-workflow-examples} illustrates this structure.
In code generation, the main agent delegates work on a CSV-export feature
to API, UI, test-authoring and documentation agents. The branches execute
in parallel, alternating model requests with file reads, edits or checks.
The main agent then integrates the changes and runs tests. In deep
research, the lead agent delegates technical, cost, security and supply
investigations, then combines the collected evidence into a
report~\cite{anthropic2025multiagent}. The number of requests and tool
calls required by each branch may vary. Therefore, after a model request
completes, its agent and the entire workflow may still be unfinished.

\paragraph{Workflow completion time.}
Let $w$ identify a workflow, $a_w$ denote its arrival time at the system,
and $c_w$ denote the time when all its required outputs are ready.
We define \emph{workflow completion time} (WCT) as the elapsed time from
arrival to completion, $\mathrm{WCT}_w=c_w-a_w$.
It includes model execution, tool execution,
queueing and dependency waits. In the coding example, integration waits
for all required changes. A delay on the last unfinished branch postpones
integration, while an earlier-finishing branch may absorb the same delay
before the join. Request latency affects WCT through these dependencies.

\begin{figure}[t]
  \centering
  \includegraphics[width=\columnwidth]{figures/agent-workflow-examples.pdf}
  \caption{Illustrative multi-agent workflows.
    Each row is one agent: $R_i$ denotes its $i$th LLM request, with
    intervening tool results carried into the next request.
    Branches have different request counts and no shared round barriers.
    The coordinator resumes after their results join.
    Positions and lengths do not encode execution time.}
  \label{fig:agent-workflow-examples}
  \Description{Two stacked left-to-right workflows. A main agent assigns
    CSV-export API, UI, test-authoring and documentation tasks in parallel,
    then integrates their outputs and runs tests. Their request counts
    are four, three, three and two, respectively. A lead researcher assigns
    technical, cost, safety and supply investigations, with four, two,
    three and two requests, before writing a report. Each branch explicitly
    unfolds its requests and intervening tools. Starts and ends are
    staggered; all edges move forward, without loops or round barriers.}
\end{figure}

\paragraph{KV reuse in multi-agent workflows.}
The model requests on these branches repeatedly consume context from
earlier steps. Coding agents carry source files, edits and test outputs
into later requests; research agents accumulate retrieved documents,
intermediate findings and tool results. Processing this growing history
from scratch on every request would repeat work already performed.
During \emph{prefill}, an LLM computes per-layer keys and values (KV)
for input tokens; \emph{decoding} extends this state as output tokens
are generated. Retaining this KV allows a subsequent request with a
matching token prefix to compute only the new suffix, provided the model
and position semantics are compatible~\cite{gao2024cachedattention,zheng2024sglang}.

This reuse occurs both within and across agents. Successive requests
from one agent can reuse KV for their shared conversation history.
Different agents can also inherit a common prompt before adding their
own tasks. For example, the API, UI and test agents in
Figure~\ref{fig:agent-workflow-examples} can share KV for common project
instructions and repository context when these form the same initial
token sequence. Different role prompts or context edits can shorten
the reusable prefix~\cite{zheng2024sglang}.
The KV cache therefore preserves computation across requests, including
while an agent executes tools or waits for another branch.

\paragraph{Failures and KV reconstruction.}
These reuse benefits depend on the cached state remaining available.
In large-scale LLM deployments, hardware faults and software crashes can
interrupt serving workers and make their GPU-resident KV
unavailable~\cite{cao2026lumen,jayakody2026ghostserve}.
For production-cluster context, Figure~\ref{fig:gpu-server-availability}
shows the public InfiniteHBD training-cluster trace: 584 recorded outages
affect 231 of 400 GPU servers over approximately 349 days, with reconstructed
availability falling to 365 servers~\cite{shou2025infinitehbd}.
The trace includes faults and operational interruptions, as detailed in
the caption; it characterizes infrastructure availability rather than
serving recovery latency.

When no reusable KV survives, a replacement worker must reconstruct it
by replaying the retained tokens, even if it already has the model weights
and the application has preserved the conversation
history~\cite{strati2024dejavu}. For a long-running agent, this means
reprocessing the context accumulated across earlier requests rather than
only the latest input~\cite{jayakody2026ghostserve}.
One failed worker can also interrupt several co-located agents, requiring
surviving workers to reconstruct multiple histories while handling normal
traffic~\cite{cao2026lumen}. The resulting delays propagate through the
workflow's dependencies: reconstruction on the last unfinished branch
postpones its join and can increase WCT.

\begin{figure}[t]
  \centering
  \includegraphics[width=\columnwidth]{figures/gpu-server-availability.pdf}
  \caption{GPU-server availability reconstructed from InfiniteHBD's
    production-training trace~\cite{shou2025infinitehbd}: 400 servers minus
    those within recorded outage intervals. Records include GPU, network,
    power/cooling, and software faults; unknown causes, tests,
    troubleshooting, and changes are also retained. Overlapping intervals
    per server are merged. Availability is measured at the node level.}
  \label{fig:gpu-server-availability}
  \Description{A step curve over approximately 349 days shows 365 to 400
    GPU servers outside active recorded fault intervals in a fixed
    400-server cohort. The dashed line marks all 400 servers.}
\end{figure}

\paragraph{KV backup across storage tiers.}
To reduce reconstruction work, a \emph{KV checkpoint} retains an already
computed prefix. If it survives the failure, recovery retrieves this
prefix and replays only the unsaved
tail~\cite{strati2024dejavu,jayakody2026ghostserve}.
KV can be retained in GPU HBM, host DRAM, or shared
storage~\cite{gao2024cachedattention,qin2025mooncake}.
These tiers differ in both storage cost and recovery latency.
HBM offers fast access but has high capacity cost and competes with
foreground inference for GPU memory. DRAM and storage-backed shared
services can provide cheaper capacity, with additional retrieval latency.

Figure~\ref{fig:motivation-recompute} illustrates the latency difference
between reconstructing a history and reusing retained KV.
At 1,010,000 tokens, full prefill takes 571.21\,s even with sparse
attention; HBM, DRAM and shared-storage reuse take 1.305, 2.663 and
10.156\,s, respectively. The history occupies 57.92\,GB across two ranks.
All paths end when the same input's KV is ready: reuse restores the
preceding history, if needed, and computes one new token.
HBM is a resident-state reference. These measurements quantify KV
readiness, excluding checkpoint creation, failure detection and routing.

\begin{figure}[t]
  \centering
  \includegraphics[width=\columnwidth]{figures/recovery-cost.pdf}
  \caption{Measured KV reconstruction and reuse costs for
    Qwen2.5-7B-Instruct-1M on two A100 PCIe 40\,GB GPUs
    (TP2, BF16, batch one), with sparse attention above 32K.
    Full prefill starts from empty KV; reuse starts from $(C-1)$ cached
    tokens and computes one. KV reconstruction time includes any
    retrieval and computation until all $C$ tokens' KV is ready.
    DRAM/shared-storage I/O is serialized with computation.
    HBM/DRAM: median and min--max of three timings;
    shared storage/full prefill: one timing per length.
    Protocol: Appendix~\ref{sec:motivation-protocol}.
    K~=~1,024; M~=~1,000,000.}
  \label{fig:motivation-recompute}
  \Description{Four measured KV reconstruction-time curves for Qwen2.5-7B-Instruct-1M
    from 8K to 1,010,000 tokens: blue circles for resident HBM reuse,
    orange squares for DRAM reuse, purple diamonds for shared-storage
    reuse, and red triangles for full prefill. Both axes are logarithmic.
    At the longest context, the four paths take 1.305, 2.663, 10.156 and
    571.21 seconds, respectively.}
\end{figure}

\subsection{Motivation and Challenges}

Limited capacity requires deciding which KV histories to retain and
which merit faster access. We use agent execution traces and measured
KV readiness costs to examine the workflow consequences of these choices.
The traces characterize branch timing; the model measurements quantify
reconstruction and reuse under a common endpoint. Together they motivate
two challenges in allocating protection resources.

\begin{figure*}[t]
  \centering
  \captionsetup[subfigure]{font=normalsize,labelfont=normalfont,
    textfont=normalfont,justification=centering,singlelinecheck=false}
  \begin{subfigure}[t]{2.48in}
    \vspace{0pt}
    \begin{minipage}[c][1.65in][c]{\linewidth}
      \centering
      \includegraphics[width=\linewidth]{figures/workflow.pdf}
    \end{minipage}
    \caption{Four-branch workflow}
    \label{fig:motivation-workflow}
  \end{subfigure}\hfill
  \begin{subfigure}[t]{1.86in}
    \vspace{0pt}
    \begin{minipage}[c][1.65in][c]{\linewidth}
      \centering
      \includegraphics[width=\linewidth]{figures/workflow-gaps.pdf}
    \end{minipage}
    \caption{Workflow-maximum gaps}
    \label{fig:motivation-branch-gap}
  \end{subfigure}\hfill
  \begin{subfigure}[t]{2.52in}
    \vspace{0pt}
    \begin{minipage}[c][1.65in][c]{\linewidth}
      \centering
      \includegraphics[width=\linewidth]{figures/recovery-exposure.pdf}
    \end{minipage}
    \caption{Recovery cost and WCT}
    \label{fig:motivation-recovery-exposure}
  \end{subfigure}
  \caption{Workflow timing and conditional recovery impact.
    (a) Four trace-derived branches with schematic Main stages A0/A5.
    Bars preserve timing; dashed tails assume an all-branch join.
    (b) ECDF of maximum branch-duration gaps per workflow: 121 workflows
    with temporally inferred fan-in; these are not failure-time slack.
    (c) Measured Prefill costs at the four branches' trace input-size
    proxies, using Figure~\ref{fig:motivation-recompute}'s configuration.
    Bars/whiskers show median/min--max over three timings after warmup.
    Hatching denotes absorbed delay; solid fill is the calculated WCT
    increase under an all-branch join with other execution fixed.
    Each case is a single-branch counterfactual, not a fault-injected WCT
    measurement. Protocol: Appendix~\ref{sec:motivation-protocol}.}

  \label{fig:motivation}
  \Description{Left: a four-branch workflow with A2 completing last at
    151.6 seconds. Middle: an ECDF of workflow-maximum branch-duration
    differences with a median of 138 seconds and a 90th percentile of
    1173 seconds. Right: A1 through A4 have measured median Prefill costs
    of 33.2, 32.6, 35.4 and 39.9 seconds. Applying each cost separately
    as an added branch pause gives calculated WCT increases of 0.0, 32.6,
    16.3 and 18.0 seconds. Hatching identifies the absorbed portion;
    whiskers show timing minima and maxima, not confidence intervals.}
\end{figure*}

\paragraph{Unequal progress across parallel branches.}
\label{sec:motivation-workflow}
Figure~\ref{fig:motivation}(\subref{fig:motivation-workflow}) shows four
trace-derived branches. A2 finishes last, 151.6\,s after the fork;
A1 finishes at 60.4\,s. If the following Main stage needs all four outputs,
A1 can absorb 91.2\,s of additional delay before postponing the join,
whereas any delay on A2 postpones it. We call the additional delay that
a branch can absorb without postponing the required output its
\emph{slack}.

The public SemiAnalysis CC/Weka trace contains 393 filtered Claude Code
sessions~\cite{semianalysis2026cctraces}. We retain overlapping subagent
groups whose branches complete and are followed by a Main request after
the last branch finishes; Main activity during the branches is allowed.
This temporal fan-in proxy retains 286 groups in 121 workflows.
For each group, we compute the longest minus the shortest branch duration,
then take each workflow's maximum gap.
Figure~\ref{fig:motivation}(\subref{fig:motivation-branch-gap}) shows a
median of 138\,s and a 90th percentile of 1,173\,s. These statistics show
substantial duration differences within the retained workflows.
The trace does not record explicit result-consumption edges, so they
neither verify all-branch joins nor measure slack at a fault. The
four-branch example is one retained group, not a representative
workflow-maximum gap.

\paragraph{From reconstruction time to workflow delay.}
To quantify the consequence of a join, we measure full prefill at the
four active requests' input-size proxies. Both this experiment and the
tier comparison in Figure~\ref{fig:motivation-recompute}
use Qwen2.5-7B-Instruct-1M~\cite{qwen2025qwen25oneM} on two A100 PCIe
40\,GB GPUs, with TP2, BF16 and batch one. Our instrumented SGLang
0.5.10rc0 execution uses MInference vertical/slash
attention~\cite{jiang2024minference} and dual chunk attention (DCA)
position handling, with sparse attention above 32K.
The four history-length proxies come from trace hash-block counts times
64 and range from 100,864 to 113,664 tokens. Each measured input is a
fixed synthetic history plus one token. After one complete warmup per
length, we collect three timings and report their median and min--max.
Appendix~\ref{sec:motivation-protocol} gives the complete protocol.

Figure~\ref{fig:motivation}(\subref{fig:motivation-recovery-exposure}) maps these costs to the
assumed all-branch join. At 29.5\,s after the fork, each case separately
adds a measured reconstruction pause $R_i$ to branch $i$, holding all
other execution and queues fixed. With offline slack $\sigma_i$ from
the trace completion times, the WCT increment is
\[
  \Delta\mathrm{WCT}_i=\max(0,R_i-\sigma_i).
\]
A1's 33.2\,s reconstruction fits within its 91.2\,s slack and leaves WCT
unchanged. A2's slightly smaller cost, 32.6\,s, is fully exposed because
its slack is zero. A3 and A4 expose 16.3 and 18.0\,s, respectively.
Thus similar reconstruction costs can have different workflow effects.
The costs are physical measurements; the WCT increments are conditional
calculations using synthetic inputs and assumed dependencies, not
fault-injected measurements of the original workflows.

\paragraph{Limitation 1: Request-level priorities incompletely express
workflow benefit.}
Recency, proximity to the next use and reconstruction cost describe when
state may be needed and how expensive it is to recover. Taken alone,
they do not determine how restoring a request sooner advances the
workflow's required outputs. An expensive history can belong to a branch
with time to spare, while a shorter recovery can delay the branch that
gates completion. Protection priorities must account for dependencies
and the branches' remaining progress.

\paragraph{Challenge 1: Estimating the workflow impact of recovery.}
A protection decision must connect the delay it avoids to the required
workflow outputs. The example uses completed traces to obtain slack,
but an online system sees only execution so far: remaining model and
tool durations are uncertain, and later dependencies may not yet be
exposed. Priorities must change as branches progress. A worker failure
can also interrupt several branches whose recoveries share resources,
so their isolated WCT increments cannot simply be added. The challenge
is to derive useful recovery priorities from partial workflow progress
while accounting for these uncertainties and interactions.

\paragraph{The cost of retaining and accessing long histories.}
\label{sec:motivation-long-context}
The tier costs in Figure~\ref{fig:motivation-recompute} distinguish
the value of retaining reusable state from the additional value of
making it quickly accessible. Retaining the history avoids most
reconstruction work; its location determines the additional retrieval
delay. At the longest context, shared storage adds 7.49\,s over DRAM
in the measured run. Whether reducing this delay advances workflow
completion depends on the recovering branch's position and progress.
The measurements characterize fixed inputs and the stated access paths;
they do not measure the net serving benefit of a protection policy.

\paragraph{Limitation 2: Recovery speed alone omits retention value and
resource costs.}
A branch that can tolerate recovery delay may still reuse its history
on the next normal request. Removing its last KV copy forfeits that reuse;
moving the copy to a slower tier preserves it but adds retrieval work.
Conversely, keeping every recoverable history in fast memory consumes
space needed by other histories and foreground execution. Minimizing
recovery latency in isolation therefore does not resolve which state
to retain, where to retain it, or how much protection the serving system
can afford.

\paragraph{Challenge 2: Balancing fast recovery, reuse and resource use.}
Finite capacity prevents every long history from remaining near
execution. Protection must decide both which histories retain reusable
state and which state merits faster access. These decisions
need to combine the workflow impact of recovery with future normal reuse,
capacity occupancy and data movement. In particular, low recovery
urgency does not imply low reuse value, and a faster restore path is
useful only to the extent that it advances required work. The challenge
is to allocate protection resources as these demands change without
allowing protection itself to undermine foreground service.

\label{sec:motivation-goals}
AegisKV targets these two challenges: using workflow progress to guide
recovery priorities, and managing limited KV capacity across storage
tiers while preserving normal reuse. The next section develops the
observations and mechanisms that support these decisions.
